\chapter{Исследовательская часть}

\section{Технические характеристики}

Технические характеристики компьютера, на котором проводилось исследование:
\begin{enumerate}[label*=---]
	\item процессор Intel Core i5-12900K;
	\item 16 Гбайт оперативной памяти;
	\item операционная система Ubuntu 22.04.
\end{enumerate}

На протежянии всего периода исследования компьютер был подключен к сети питания и количество фоновых процессов было сведено к минимуму.

\section{Методика исследования}

Для оценки эффективности разработанного метода обеспечение дифференциальной приватности при генерации синтетических данных.
Исследование включало в себя измерение следующих характеристик:
\begin{enumerate}[label*=---]
	\item время работы метода при различных размерах данных (количества записей и количества атрибутов);
	\item качество генерируемых синтетических данных;
\end{enumerate}
Также исследование включало сравнение с альтернативными методами генерации без дифференциальной приватности.

При проведении исследования:
\begin{enumerate}[label*=---]
	\item количество записей в исходных данных было равно 10000, за исключением тех замеров, в которых оно изменялось для исследования зависимостей;
	\item количество атрибутов в исходных данных было равно 9, за исключением тех замеров, в которых оно изменялось для исследования зависимостей;
	\item параметр приватности был равен 1, за исключением тех замеров, в которых он изменялся для исследования зависимостей;
	\item cтепень Байесовской сети был равен 2, за исключением тех замеров, в которых она изменялась для исследования зависимостей;
	\item число генерируемых записей было равно 100000, за исключением тех замеров, в которых оно изменялось для исследования зависимостей;
\end{enumerate}

%	\item распределения признаков (гистограммы, средние, дисперсии) между реальными и синтетическими данными.
	

%Сравнение корреляций:
%Постройте корреляционные матрицы для реальных и синтетических данных, оцените сохранение взаимосвязей.
%Сравнение многомерных зависимостей:
%Используйте методы визуализации (например, PCA, t-SNE) для сравнения структуры данных.	

%Влияние параметра ε (эпсилон):
%Проведите серию экспериментов с разными значениями ε (например, 0.01, 0.1, 1, 10) и оцените, как меняется качество синтетических данных и уровень приватности.
%Влияние порога категориальности (threshold):
%Изучите, как изменение порога влияет на качество моделирования категориальных признаков.

%Влияние степени сети (degree_of_bayesian_network):
%Изучите, как максимальное число родителей в сети влияет на сохранение зависимостей и качество синтетических данных.

%Атаки восстановления (Membership Inference Attack):
%Проведите эксперименты по восстановлению исходных данных из синтетических, чтобы оценить уровень приватности.
%Сравнение с другими методами синтеза:
%Сравните ваш подход с другими генераторами синтетических данных (например, CTGAN, Synthpop и др.) по качеству и приватности.



\section{Исследование характеристик}

\subsection{Время работы}

Одной из ключевых характеристик реализованного метода является зависимость время работы от разных входных параметров.

Для исследования данной зависимости были проведены измерения времени работы для:
\begin{enumerate}
	\item исходных датасетов разного размера количества записей;
	\item исходных датасетов разного размера количества атрибутов;
	\item разных параметрах приватности;
	\item разных степенях БС.
\end{enumerate}
\newpage
На рисунке~\ref{img:new-r-size} представлена диаграмма зависимости времени работы от размера количества записей входного датасета:

\includeimage
{new-r-size} % Имя файла без расширения (файл должен быть расположен в директории inc/img/)
{f} % Обтекание (с обтеканием)
{h!} % Положение рисунка (см. wrapfigure из пакета wrapfig)
{0.8\textwidth} % Ширина рисунка
{Зависимость времени работы от размера количества записей входного датасета} % Подпись рисунка
% \imgs{index-struct}{h!}{1}{Индекс как структура данных}

На рисунке~\ref{img:new-r-atr} представлена диаграмма зависимости времени работы программы от размера количества атрибутов входного датасета:

\includeimage
{new-r-atr} % Имя файла без расширения (файл должен быть расположен в директории inc/img/)
{f} % Обтекание (с обтеканием)
{h!} % Положение рисунка (см. wrapfigure из пакета wrapfig)
{0.8\textwidth} % Ширина рисунка
{Зависимость времени работы программы от размера количества атрибутов входного датасета} % Подпись рисунка
% \imgs{index-struct}{h!}{1}{Индекс как структура данных}

\newpage
Результаты измерений показывают, что: время работы увеличивается вместе с количеством атрибутов в датасета. Также из двух диаграмм видно, что при увеличении количества атрибутов на выполнение работы уходит больше времени, чем при увеличении количества записей.

\newpage
На рисунке~\ref{img:new-r-eps} представлена диаграмма зависимости времени работы от разных параметров приватности при одинаковых других входных параметрах:

\includeimage
{new-r-eps} % Имя файла без расширения (файл должен быть расположен в директории inc/img/)
{f} % Обтекание (с обтеканием)
{h!} % Положение рисунка (см. wrapfigure из пакета wrapfig)
{0.8\textwidth} % Ширина рисунка
{Зависимость времени работы от разных значений параметра приватности}

На рисунке~\ref{img:new-r-degree} представлена диаграмма зависимости времени работы от степени Байесовской сети:

\includeimage
{new-r-degree} % Имя файла без расширения (файл должен быть расположен в директории inc/img/)
{f} % Обтекание (с обтеканием)
{h!} % Положение рисунка (см. wrapfigure из пакета wrapfig)
{0.8\textwidth} % Ширина рисунка
{Зависимость времени работы от степени Байесовской сети}

\newpage
\subsection{Качество данных}

%Для оценки качества генерируемых синтетических данных используются следующие метрики: средняя абсолютная ошибка (MAE) между распределениями исходных и синтетических данных и коэффициент детерминации $R^2$.


%Cредняя абсолютная ошибка (англ. mean absolute error сокр. MAE) --- оценка качества результата, которая оценивает среднее отклонение синтетических корреляций от реальных.
%Коэффициент детерминации (сокр. $R^2$) --- это статистическая мера, которая показывает, какую долю дисперсии зависимой переменной объясняет модель по сравнению с простым предсказанием средним значением.~\cite{mae}


%Средняя абсолютная ошибка вычисляется по формуле~\ref{eq:mae}:
%\begin{equation}\label{eq:mae}
	%MAE = \frac{1}{n}\sum_{i=1}^{n}|y_i - \hat{y}_i|
%\end{equation}

%Коэффициент детерминации $R^2$ вычисляется по формуле~\ref{eq:r2}:
%%R^2 = 1 - \frac{\sum_{i=1}^{n}(y_i - \hat{y}_i)^2}{\sum_{i=1}^{n}(y_i - \bar{y})^2}
%\end{equation}
%где $y_i$ --- исходные значения; $\hat{y}_i$ --- синтетические значения; $\bar{y}$ --- среднее значение исходных данных.


Исследуется зависимость качества генерируемых синтетических данных от значения параметра приватности $\varepsilon$. Результаты измерений представлены в таблице~\ref{tab:varepsilon_results} и на рисунке~\ref{img:new-r2-mae}.
\begin{table}[h!]
	\centering
	\captionsetup{justification=raggedleft}
	\caption{Влияние параметра $\varepsilon$ на качество синтетических данных}
	\label{tab:varepsilon_results}
	\begin{tabular}{|c|c|}
		\hline
		Параметр $\varepsilon$  & Фробениусова норма  \\
		\hline
		0.1 &  0.31 \\
		\hline
		0.5 & 0.22  \\
		\hline
		1.0 & 0.16  \\
		\hline
		2.0 & 0.12  \\
		\hline
		5.0  & 0.07 \\
		\hline
		10.0  & 0.03  \\
		\hline
	\end{tabular}
\end{table}
\includeimage
{new-r2-mae} % Имя файла без расширения
{f} % Обтекание
{h!} % Положение рисунка
{0.9\textwidth} % Ширина рисунка
{Влияние параметра $\varepsilon$ на качество синтетических данных} % Подпись рисунка


%Результаты эксперимента показывают обратную зависимость между уровнем приватности и качеством данных:
%\begin{enumerate}[label*=---]
	%\item При $\varepsilon = 0.1$ (высокая приватность): $R^2 = 0.42$, $MAE = 0.31$;
	%\item При $\varepsilon = 1.0$ (средняя приватность): $R^2 = 0.74$, $MAE = 0.16$;
	%\item При $\varepsilon = 10.0$ (низкая приватность): $R^2 = 0.92$, $MAE = 0.07$.
%\end{enumerate}

Оптимальное значение $\varepsilon = 1.0{-}2.0$ обеспечивает баланс между приватностью и полезностью данных.

% \subsection{Сохранение корреляционной структуры}
% Для оценки сохранения корреляционной структуры используется Фробениусова норма, которая позволяет количественно оценить, насколько одна матрица корреляций отличается от другой.
% Фробениусова норма разности корреляционных матриц определяется выражением~\ref{eq:frob}:
% \begin{equation}\label{eq:frob}
	%||\Delta R||_F = \sqrt{\sum_{i,j}(R_{orig}(i,j) - R_{synth}(i,j))^2}
%\end{equation}
%где $R_{orig}$ --- корреляционная матрица исходных данных; $R_{synth}$ --- корреляционная матрица синтетических данных.


%Результаты исследования представлены в таблице~\ref{tab:correlation} и на рисунке~\ref{img:frob}.

%\begin{table}[h!]
	%\centering
	%\caption{Сохранение корреляционной структуры}
	%\label{tab:correlation}
	%\begin{tabular}{|c|c|}
		%\hline
	%Параметр $\varepsilon$ & Фробениусова норма \\
	%\hline
	%0.1 & 0.87 \\
	%\hline
	%1.0 & 0.34 \\
	%\hline
	%10.0 & 0.12 \\
	%\hline
%\end{tabular}
%\end{table}

%\includeimage
%{frob} % Имя файла без расширения
%{f} % Обтекание
%{h!} % Положение рисунка
%{0.8\textwidth} % Ширина рисунка
%{Сохранение корреляционной структуры при различных уровнях параметраприватности} % Подпись рисунка

%Метод успешно сохраняет основные корреляционные зависимости даже при высоком уровне приватности. При $\varepsilon = 1.0$ фробениусова норма составляет 0.34, что указывает на хорошее сохранение структуры данных.

%На рисунке~\ref{img:r-degree} представлена диаграмма зависимости MAE от степени Байесовской сети:

%Для исследования данной зависимости были проведены измерения для:

%\includeimage
%{mae_vs_degree} % Имя файла без расширения (файл должен быть расположен в директории inc/img/)
%{f} % Обтекание (с обтеканием)
%{h!} % Положение рисунка (см. wrapfigure из пакета wrapfig)
%{0.8\textwidth} % Ширина рисунка
%{Зависимость MAE от степени Байесовской сети}

\section{Сравнение с альтернативными методами}

Проводится сравнение разработанного метода с альтернативами:
\begin{enumerate}
	\item ctGAN без дифференциальной приватности;
	\item Предложенный метод с $\varepsilon = 1.0$.
	\item VAE без дифференциальной приватности;
\end{enumerate}
\newpage
Результаты сравнения представлены в таблице~\ref{tab:baseline_comparison} и на рисунке~\ref{img:new-meth1} и~\ref{img:new-meth2}:

\begin{table}[h!]
	\centering
	\captionsetup{justification=raggedleft}
	\caption{Сравнение с альтернативными методами}
	\label{tab:baseline_comparison}
	\begin{tabular}{|c|c|c|}
		\hline
		Метод & Фробениусова норма & Время генерации, мс \\
		\hline
		Исходные данные & 0.00 & --- \\
		\hline
		ctGAN без ДП & 0.08 & 45211 \\
		\hline
		Предложенный метод & 0.12 & 542437 \\
		\hline
		VAE без ДП & 0.09 & 50109 \\
		\hline

	\end{tabular}
\end{table}

\includeimage
{new-meth1} % Имя файла без расширения
{f} % Обтекание
{h!} % Положение рисунка
{0.8\textwidth} % Ширина рисунка
{Сравнение разработанного метода с альтернативами} % Подпись рисунка
\includeimage
{new-meth2} % Имя файла без расширения
{f} % Обтекание
{h!} % Положение рисунка
{0.8\textwidth} % Ширина рисунка
{Сравнение разработанного метода с альтернативами} % Подпись рисунка


\newpage
\section{Выводы}

В данном разделе были проведены исследования характеристик разработанного метода обеспечения дифференциальной приватности при генерации синтетических данных.
Основные результаты исследования:
\begin{enumerate}[label*=---]
	\item Результаты измерений показывают, что: время работы увеличивается вместе с количеством записей в датасет
	\item Результаты измерений показывают, что: время работы увеличивается вместе с количеством атрибутов в датасета
	\item Оптимальное значение $\varepsilon = 1.0{-}2.0$ обеспечивает баланс между приватностью и полезностью данных;
	% \item Метод успешно сохраняет основные корреляционные зависимости даже при высоком уровне приватности. При $\varepsilon = 1.0$ фробениусова норма составляет 0.34, что указывает на хорошее сохранение структуры данных;
	\item Разработанный метод показывает высокое качество данных при обеспечении дифференциальной приватности. Точность составляет 0.12, что является приемлемым компромиссом.
\end{enumerate}

